\message{ !name(main.tex)}%%%%%%%%%%%%%%%%%%%%%%%%%%%%Document Info%%%%%%%%%%%%%%%%%%%%%%%%%%%%%%%%%%%%%%%
\documentclass[a4paper, 11pt]{report}
\input{~/LaTeX-Template/headers/note.tex}
\usepackage{animate}

%%%%%%%%%%%%%%%%%%%%%%%%%%%%Article Info%%%%%%%%%%%%%%%%%%%%%%%%%%%%%%%%%%%%%%%%
\author{Bhashkar Paudyal}
\title{Calculus II}

%%%%%%%%%%%%%%%%%%%%%%%%%%%%Misc%%%%%%%%%%%%%%%%%%%%%%%%%%%%%%%%%%%%%%%%%%%%%%%%
\thispagestyle{empty}
\addbibresource{reference.bib}

%%%%%%%%%%%%%%%%%%%%%%%%%%%%Additional Packages%%%%%%%%%%%%%%%%%%%%%%%%%%%%%%%%%
\usepackage{adjustbox}
\usepackage{centernot}


%%%%%%%%%%%%%%%%%%%%%%%%%%%%Document Begin%%%%%%%%%%%%%%%%%%%%%%%%%%%%%%%%%%%%%%
\begin{document}

\message{ !name(lec_5.tex) !offset(-22) }
\chapter{Line, Circle, Cylinder and Sphere}
\lecture{5}{9 May. 00:51}{Fifth Lecture}
\label{chap:something}

\section{Straight Lines in Polar Form}
\label{sec:straight-lines-

In Cartesian Coordinate system, a straight line can be representation in many
form. One of them is slope form \( y = mx + c \). But the one that mostly
relates with Polar System is the normal form of representing a straight line.
\[ x \cos(\alpha) + y \cos(\alpha) \]

\message{ !name(main.tex) !offset(9) }

\end{document}
